\documentclass[10pt]{article}
% Estilo compartilhado pelas versoes PT e EN
\usepackage[a4paper,top=1.3cm,bottom=1.3cm,left=1.5cm,right=1.5cm]{geometry}
\usepackage[T1]{fontenc}
\usepackage[utf8]{inputenc}
\usepackage{lmodern}
\usepackage{titlesec}
\usepackage{enumitem}
\usepackage{xcolor}
\usepackage[hidelinks]{hyperref}
\usepackage{tabularx}

\definecolor{linkblue}{RGB}{20,70,150}
\newcommand{\link}[2]{\href{#1}{\textcolor{linkblue}{#2}}}

\pagestyle{empty}
\setlength{\parindent}{0pt}
\linespread{1.04}

% Titulos de secao em caixa alta com linha abaixo
\titleformat{\section}{\large\bfseries\raggedright}{}{0em}{\MakeUppercase}[\vspace{1pt}\titlerule]
\titlespacing*{\section}{0pt}{10pt}{4pt}

% \entrada{titulo em negrito}{detalhe}{data a direita}
\newcommand{\entrada}[3]{%
  \par\vspace{4pt}\noindent\textbf{#1}#2\hfill#3\par\vspace{1pt}}

% Itens com ponto medio, como no modelo
\newenvironment{itens}{\begin{itemize}[leftmargin=0pt,label={$\cdot$},itemsep=1pt,topsep=1pt,parsep=0pt,itemindent=8pt,labelsep=4pt]}{\end{itemize}}

% Linha "Rotulo   conteudo" para competencias
\newcommand{\linha}[2]{\noindent\begin{tabularx}{\linewidth}{@{}p{4.0cm}@{}X@{}}\textbf{#1} & #2\end{tabularx}\par\vspace{0pt}}

% hifenizacao em portugues indisponivel aqui; no Overleaf use \usepackage[brazilian]{babel}

\begin{document}

\begin{center}
  {\LARGE\bfseries MATHEUS GOMES FÉLIX DA SILVA}\\[5pt]
  +55 83 99106-0871 \,-\, \link{mailto:matheusgfsilva16@gmail.com}{matheusgfsilva16@gmail.com} \,-\, João Pessoa, PB
\end{center}
\vspace{2pt}

\linha{Idiomas}{Português (nativo), Inglês (intermediário avançado).}
\textbf{GitHub}: \link{https://github.com/Matheusgfsilva}{Matheusgfsilva} \quad
\textbf{LinkedIn}: \link{https://www.linkedin.com/in/matheus-gomesfs/}{in/matheus-gomesfs} \quad
\textbf{Site}: \link{https://matheusgfsilva.github.io}{matheusgfsilva.github.io}

\section{Resumo}
Estudante de Ciência da Computação na UFPB e Machine Learning Engineer voluntário no laboratório ARIA, com foco em visão computacional (detecção, segmentação e OCR) e modelos preditivos. Experiência full stack como sócio desenvolvedor do Compendium. Comunicativo, disciplinado e com foco em empreender em tecnologia.

\section{Formação}
\entrada{Bacharelado em Ciência da Computação}{ — Universidade Federal da Paraíba (UFPB)}{nov. 2025 – jun. 2029}

\section{Experiência}
\entrada{Machine Learning Engineer (Voluntário) - Laboratório ARIA}{ [Remoto]}{mai. 2026 – atual}
\begin{itens}
  \item Desenvolvo modelos de visão computacional para leitura automática de nomes e dosagens de medicamentos em imagens de seringas, bombas de infusão, bolsas, frascos e canetas, em parceria com o Hospital Israelita Albert Einstein.
  \item Treino e comparo modelos de detecção e segmentação (YOLO, RF-DETR) para identificar objetos e líquidos em ambiente hospitalar.
  \item Integro OCR ao pipeline de detecção e conduzo análises de desempenho em benchmarks para comparar arquiteturas e orientar a escolha do modelo.
  \item Desenvolvo modelos de análise e previsão de dados financeiros, como nível de confiança de clientes bancários e limite de crédito esperado.
\end{itens}

\entrada{Sócio e Desenvolvedor - \link{https://usecompendium.com}{Compendium}}{}{mar. 2026 – set. 2026}
\begin{itens}
  \item Cofundei e desenvolvi uma aplicação web multi-tenant (um subdomínio por cliente) que centraliza normas, procedimentos, reuniões, advertências e suspensões.
  \item Implementei o controle de leitura das normas por assinatura eletrônica com validade jurídica (Autentique), com os PDFs assinados no Cloudflare R2.
  \item Desenvolvi a gestão de cargos e permissões, que define o que cada cargo pode visualizar e editar no sistema.
  \item Integrei cobrança (Asaas), envio de e-mails (Resend) e notificações push em PWA, sobre uma API REST com autenticação JWT.
  \item \textbf{Stack:} JavaScript, Node.js (Fastify), React (Vite, Tailwind CSS), PostgreSQL, Prisma, Railway.
\end{itens}

\section{Projetos}
\entrada{\link{https://github.com/Matheusgfsilva/CONTROLE_FINANCEIRO}{Controle Financeiro}}{}{Python, Rich}
\begin{itens}
  \item Sistema de gestão financeira para pequenos negócios, com receitas, despesas, saldo, relatórios e categorias editáveis.
  \item Arquitetura modular com cerca de 3.000 linhas, validação de entradas, persistência em JSON e geração de dados de teste.
\end{itens}

\entrada{\link{https://github.com/Matheusgfsilva/CONTROLE_CLIENTES}{Controle de Clientes}}{}{Python, Rich}
\begin{itens}
  \item Sistema CRUD de cadastro de clientes, com validação de CPF, celular e e-mail e bloqueio de CPF duplicado.
\end{itens}

\section{Competências}
\linha{Linguagens}{Python, JavaScript, Java, SQL.}
\linha{Machine Learning}{Visão computacional, detecção e segmentação, OCR, YOLO, RF-DETR, benchmarking, modelos preditivos.}
\linha{Desenvolvimento web}{Node.js, Fastify, React, Tailwind CSS, APIs REST, PostgreSQL, MySQL, Prisma.}
\linha{Automação}{n8n (nível inicial): criação de fluxos e integração entre serviços.}
\linha{Comportamentais}{Comunicação, relacionamento interpessoal, disciplina e autogestão.}

\end{document}
